\section{Introduction}
\label{sec:intro}

The method of fundamental solutions (MFS) approximates the solution of a boundary value problem by a sum of point sources placed outside the domain \cite{kupradze1964,fairweather1998}. Each source satisfies the differential equation exactly, so only the boundary condition has to be enforced, at a set of points on the boundary. There is no mesh and no integration. For scattering problems the sources lie inside the scatterer and radiate into the exterior \cite{fairweather2003}. In electromagnetics the same idea is known as the method of auxiliary sources (MAS) or the method of discrete sources \cite{kaklamani2002,doicu2000}.

When the MFS works it converges very quickly. For the Helmholtz equation on analytic domains, Barnett and Betcke \cite{barnett2008} proved exponential convergence for the disc and explained when the method is stable. The solution of the boundary value problem can be continued analytically beyond the boundary, up to a set of singularities. If the sources enclose these singularities, the coefficients stay bounded and the error decays exponentially. If they do not, the coefficients grow exponentially and accuracy is lost in floating-point arithmetic. The MAS literature states the same rule: the auxiliary surface must enclose the singularities of the continued scattered field \cite{leviatan1990,zaridze2002,fikioris2006,tsitsas2018,kyurkchan1996}.

This rule has a consequence that is easy to state. A body with an edge has singularities of the continued field on the edge itself, so no surface inside the body can enclose them. The usual MFS layout, with sources on a shrunk copy of the surface, cannot converge. The MAS literature avoids the difficulty by rounding the edges \cite{bogdanov1999}. In two dimensions the difficulty has been resolved. Lightning solvers place sources in clusters that approach each corner exponentially and reach root-exponential convergence \cite{gopal2019siam,gopal2019pnas}, also for the Helmholtz equation \cite{ginn2023}; Hochman, Leviatan and White \cite{hochman2013} obtained similar layouts earlier by rational fitting. In three dimensions the question is open. Gopal and Trefethen close their paper with the sentence ``The extension to 3D is a bigger issue to be considered'' \cite{gopal2019siam}. Brubeck and Trefethen write that it is ``undoubtedly possible, but not straightforward'' \cite{brubeck2022}. Kong and Shen \cite{kong2026} treat isolated point singularities in three dimensions and call edges and vertices of polyhedra ``particularly challenging''.

For Maxwell's equations the MFS has been used on smooth bodies. Recent examples are the regularised formulation of Ye et al.\ \cite{ye2023} and the solver of Weed et al.\ \cite{weed2026} for layered periodic media, and Kangro \cite{kangro2012} gives a convergence result for analytic surfaces. We are not aware of an MFS for Maxwell scattering that converges on a body with edges or tips.

This paper studies the MFS for the scattering of a plane wave by a perfect electric conductor in three dimensions. It grew out of the first author's MSci thesis \cite{thesis2025}, which applied the MFS to a sphere, a cube and the NASA almond \cite{woo1993}. We keep these shapes and add a torus. The four shapes form a sequence in which the singularities of the continued field move from the centre of the body to its surface:
\begin{itemize}
  \item \emph{Sphere} (\autoref{sec:sphere}). The only singularity is the centre. Dipoles on any inner sphere converge like $\exp(-c\sqrt{N})$ in the number of sources $N$, down to an error of $2\times10^{-14}$ against the Mie series.
  \item \emph{Torus} (\autoref{sec:torus}). The singularities fill a ring inside the tube. Dipoles on an inner torus converge to $7\times 10^{-12}$, provided that its radius exceeds a threshold which we give in closed form.
  \item \emph{Almond} (\autoref{sec:almond}). The singularities lie on a curve in the mid-plane, within 6\% of the rim, and reach the surface at the tip. No inner copy of the surface works. Sources on the mid-plane, clustered towards this curve, do.
  \item \emph{Cube} (\autoref{sec:cube}). The singularities reach the surface along all twelve edges. Sources on the bisector planes of the edges, clustered towards the edges, converge.
\end{itemize}

Our contributions are the following.
\begin{enumerate}
  \item We show that the scattered field of a convex conducting polyhedron is exactly the field of electric and magnetic currents on its interior medial surface (\autoref{thm:skeleton}). This explains where the sources must be placed, why a shrunk surface fails, and why both electric and magnetic dipoles are needed.
  \item We introduce the \emph{singular tail}. Sources that cluster towards an edge or a tip must stop at some smallest distance, and the cut-off limits the accuracy. Below the cut-off we continue the sources, but tie them together with the power laws of the local field expansion, so that they add only a few unknowns. The exponents are not fitted. For an edge they follow from the Meixner expansion \cite{meixner1972}, and for the cube vertex and the almond tip we compute them from the eigenvalues of the Laplace--Beltrami operator outside the corresponding cone \cite{desmedt1986}.
  \item With these layouts the MFS converges on the cube to a relative error of about $10^{-4}$ in the scattered field, and it agrees to $2\times 10^{-4}$ with a boundary element solution on an edge-graded mesh. On the almond the computed monostatic radar cross-section agrees with a published reference simulation \cite{kelley2020} to between 0.01 and 0.14\,dB on average at four frequencies up to 10.25\,GHz.
  \item We compare accuracy against computing time with the boundary element library Bempp \cite{betcke2021} on all four shapes, and we state plainly where the MFS is ahead and where it is not.
\end{enumerate}

\autoref{sec:formulation} states the problem and the method. \autoref{sec:skeleton} explains why inner surfaces fail for the cube and the almond. \autoref{sec:discussion} lists the limits of what we have shown.
